\documentclass[11pt]{article}
\usepackage[margin=1in]{geometry}
\usepackage{amsmath,amssymb,amsthm}
\usepackage{hyperref}
\newtheorem{theorem}{Theorem}
\newtheorem{lemma}[theorem]{Lemma}
\newtheorem{proposition}[theorem]{Proposition}
\newtheorem{corollary}[theorem]{Corollary}
\theoremstyle{definition}
\newtheorem{definition}[theorem]{Definition}
\theoremstyle{remark}
\newtheorem{remark}[theorem]{Remark}

\title{Disproof of Conjecture 00000001564:\\
a $k$-dimensional cs-neighborly polytope can have only $2k<2^{k+1}$ vertices}
\author{jilint777}
\date{2026-10-04}

\begin{document}
\sloppy
\maketitle

\begin{abstract}
Conjecture 00000001564 asserts that the minimal number of vertices of a
$k$-dimensional cs-neighborly polytope is $2^{k+1}$, and that the extremal
bodies are the Hanner polytopes. The first clause is false for every $k\ge 1$.
The $k$-dimensional cross-polytope $C_k^\Delta=\operatorname{conv}\{\pm e_1,\dots,\pm e_k\}$
is centrally symmetric, has $2k$ vertices, and is cs-neighborly in the
strongest possible sense: every set of its vertices containing no antipodal pair
is the vertex set of a face. Since $2k<2^{k+1}$, the claimed minimum is wrong;
the smallest cases are the square ($4<8$) and the octahedron ($6<16$). The
counterexample works under every reading of ``cs-neighborly'' we know, including
the statement's own definition. We also show that the true minimum is $2k$, and
that no Hanner polytope has $2^{k+1}$ vertices. The refutation is formalized in
Lean~4 (core library only) for $k=1,2,3,4$. The Chinese version of the statement
asks about \emph{$2k$-dimensional} polytopes. That claim holds for $k=1,2$ and fails for
every $k\ge3$; the smallest counterexample is the $6$-dimensional cross-polytope
($12<16$), also formalized in Lean.
\end{abstract}

\section{The conjecture}
\begin{quote}
\emph{Definition:} A cs-neighborly polytope is centrally symmetric and has
support points in every direction. \emph{Conjecture:} The minimal number of
vertices of a $k$-dimensional cs-neighborly polytope is $2^{k+1}$; and the
extremal bodies are the Hanner polytopes.
\end{quote}
The Chinese version differs in one point: it speaks of a ``2$k$-dimensional''
polytope (the dimension is written as ``two $k$''), while the number of vertices is still
$2^{k+1}$. We treat this reading separately in Section~\ref{sec:cn}. The conjecture is a
conjunction, so it is enough to refute the first clause:
\begin{equation}\label{eq:claim}
\textstyle\min\{\,f_0(P) : P \text{ a $k$-dimensional cs-neighborly polytope}\,\} = 2^{k+1}
\qquad(k\ge 1),
\end{equation}
where $f_0(P)$ is the number of vertices of $P$. A statement ``the minimum is $n$''
contains the lower bound ``every such $P$ has $f_0(P)\ge n$''. We refute this lower bound,
which is the weakest part of~\eqref{eq:claim}.

\paragraph{Readings of ``cs-neighborly''.}
The definition in the statement is unusual. ``Having support points in every
direction'' holds for every nonempty compact set, so under the literal reading a
cs-neighborly polytope is simply a centrally symmetric polytope. The standard notion is
Gr\"unbaum's~\cite{Gr}. A centrally symmetric (cs) polytope $P$ is
\emph{cs-$j$-neighborly} if every set of at most $j$ vertices that contains no pair of
antipodal vertices $\{v,-v\}$ is the vertex set of a face of $P$. ``cs-neighborly'' means
cs-$2$-neighborly (any two non-antipodal vertices span an edge), or cs-$\lfloor k/2\rfloor$-neighborly
for a $k$-polytope, depending on the author. We treat all these readings at
once, via the strongest property:
\begin{itemize}
\item[(F)] every nonempty set of vertices without an antipodal pair is the vertex set of a face.
\end{itemize}
Property (F) implies cs-$j$-neighborliness for every $j$, and it implies the
statement's literal definition. The cross-polytope has property (F)
(Theorem~\ref{thm:cross}), so it is a counterexample under every reading.

\section{Definitions}
A \emph{polytope} is the convex hull $P=\operatorname{conv}V$ of a finite set
$V\subset\mathbb R^k$. Its \emph{dimension} is the dimension of its affine hull. A
point $v\in V$ is a \emph{vertex} of $P$ if some linear functional $c$ attains its
maximum over $P$ only at $v$. Every vertex of $P$ lies in $V$. If $V$ is in
\emph{convex position} (every point of $V$ is a vertex), then $f_0(P)=|V|$. For a
functional $c$, the set $F_c=\{x\in P: c\cdot x=\max_P c\}$ is a face of $P$. Its
vertex set is $F_c\cap V=\operatorname{argmax}_{v\in V} c\cdot v$, and every
nonempty face of a polytope arises this way. So a set $S\subseteq V$ is the vertex set of a
nonempty face iff $S=\operatorname{argmax}_{V}c$ for some $c$. $P$ is \emph{centrally
symmetric} if $P-x_0=-(P-x_0)$ for some center $x_0$. For $x_0=0$ this means $V=-V$
for the vertex set $V$.

\section{The cross-polytope}
\begin{theorem}\label{thm:cross}
Let $k\ge1$ and $V_k=\{\pm e_1,\dots,\pm e_k\}\subset\mathbb R^k$, $C^\Delta_k=\operatorname{conv}V_k$.
Then:
\begin{enumerate}
\item $C^\Delta_k$ is $k$-dimensional and centrally symmetric (about $0$);
\item $V_k$ is in convex position, so $f_0(C^\Delta_k)=2k$;
\item every nonempty $S\subseteq V_k$ containing no antipodal pair is the vertex
set of a face, namely $S=\operatorname{argmax}_{V_k}c_S$ with $c_S=\sum_{s\in S}s$;
\item $C^\Delta_k$ has a support point in every direction.
\end{enumerate}
\end{theorem}
\begin{proof}
(1) $V_k=-V_k$. The affine hull contains $0=\tfrac12(e_1+(-e_1))$ and every $e_i$,
so it is $\mathbb R^k$.
(2) For $v\in V_k$ we have $v\cdot v=1$ and $v\cdot w\in\{0,-1\}$ for $w\in V_k\setminus\{v\}$.
So $c=v$ is maximized over $V_k$ only at $v$, and hence over $C^\Delta_k$ only at $v$.
(3) Since $S$ contains at most one of $e_i,-e_i$, the $i$-th coordinate of $c_S$
is $+1$ if $e_i\in S$, $-1$ if $-e_i\in S$, and $0$ otherwise. Hence
$c_S\cdot s=1$ for $s\in S$. For $w\in V_k\setminus S$ we get $c_S\cdot w\in\{0,-1\}$:
it is $-1$ if $-w\in S$ and $0$ otherwise. So $\operatorname{argmax}_{V_k}c_S=S$.
(4) Any functional attains its maximum on the compact set $C^\Delta_k$.
\end{proof}

\section{The disproof}
\begin{theorem}\label{thm:main}
For every $k\ge1$ there is a $k$-dimensional centrally symmetric polytope with
property (F) (hence cs-$j$-neighborly for every $j$, and cs-neighborly in the sense
of the statement) that has $2k<2^{k+1}$ vertices. Hence claim~\eqref{eq:claim} is false
for every $k\ge1$, and so is the conjecture.
\end{theorem}
\begin{proof}
Take $C^\Delta_k$ (Theorem~\ref{thm:cross}). We have $2k\le 2^k<2^{k+1}$, since
$2\cdot1=2^1$ and $2(k+1)=2k+2\le 2^k+2^k$.
\end{proof}
The smallest instances are the segment ($k=1$: $2<4$), the square
$\operatorname{conv}\{(\pm1,0),(0,\pm1)\}$ ($k=2$: $4<8$; its $8$ faces with
antipodal-free vertex sets are the $4$ vertices and the $4$ edges), and the octahedron ($k=3$:
$6<16$; all $26$ antipodal-free vertex sets are faces: $6$ vertices, $12$ edges, $8$ triangles).
For $k=0$ the only polytope is a point, with $1\ne2^1$ vertex.

\begin{remark}[Other readings]\label{rem:readings}
\begin{enumerate}
\item \emph{Center.} The cross-polytope is symmetric about the origin, so it is
centrally symmetric about some center under any convention.
\item \emph{$k$ as the neighborliness parameter.} If ``$k$-dimensional'' is read as
``cs-$k$-neighborly'' in arbitrary dimension, then $C^\Delta_k$ is a cs-$k$-neighborly
polytope with $2k<2^{k+1}$ vertices.
\item \emph{Lattice vs.\ real polytopes.} Our example has integer vertices. A
counterexample in a subclass is a counterexample to the universally quantified claim.
\item \emph{Facets instead of vertices.} This is not what is written, but it also fails.
$C^\Delta_k$ has $2^k<2^{k+1}$ facets. Under the literal definition, the cube $[-1,1]^k$
qualifies and has $2k$ facets.
\end{enumerate}
\end{remark}

\section{\texorpdfstring{The Chinese reading: dimension $2k$}{The Chinese reading: dimension 2k}}\label{sec:cn}
The Chinese text reads: the minimal number of vertices of a $2k$-dimensional cs-neighborly
polytope is $2^{k+1}$. As a universal statement, this is
\begin{equation}\label{eq:cn}
\forall k\ge1:\quad \min\{\,f_0(P) : P \text{ a $2k$-dimensional cs-neighborly polytope}\,\}=2^{k+1},
\end{equation}
and in particular it contains the universal lower bound: every $2k$-dimensional cs-neighborly
polytope has at least $2^{k+1}$ vertices.

\begin{theorem}\label{thm:cn}
Claim~\eqref{eq:cn} holds for $k=1,2$ and fails for every $k\ge3$, under each reading of
``cs-neighborly''. The smallest counterexample is the $6$-dimensional cross-polytope
$C^\Delta_6$, which has property (F) and $12<16=2^{4}$ vertices.
\end{theorem}
\begin{proof}
By Proposition~\ref{prop:true} (below), the minimum in dimension $2k$ is $4k$, under every
reading. We have $4k=2^{k+1}$ for $k=1,2$ ($4=4$, $8=8$), and $4k<2^{k+1}$ for $k\ge3$
($12<16$, and if $4k<2^{k+1}$ then $4(k+1)=4k+4<2^{k+1}+2^{k+1}=2^{k+2}$, using
$4\le 2^{k+1}$). For $k\ge3$, $C^\Delta_{2k}$ is a counterexample (Theorem~\ref{thm:cross}).
\end{proof}
So the Chinese version is also false as a statement about all $k$, but its first two instances
are true; a counterexample needs dimension at least $6$. Lean proves
\texttt{conjecture\_00000001564\_false\_2k}: for every notion $N$ that $C^\Delta_6$ satisfies,
$\neg\,\forall k\ge1,\ \texttt{MinVerticesIs N (2*k) (2\^{}(k+1))}$, with instances for the
literal definition, property (F), cs-$j$-neighborliness for every $j$, and cs-$3$-neighborliness
(Gr\"unbaum's cs-$\lfloor d/2\rfloor$-neighborliness for $d=6$). The case $k=3$ needs only the
lower-bound half, which $C^\Delta_6$ violates. The truth of the instances $k=1,2$ rests on
Proposition~\ref{prop:true}, which is not formalized in dimensions $2$ and $4$.

\section{The true minimum and the Hanner clause}
\begin{proposition}\label{prop:true}
Every $k$-dimensional centrally symmetric polytope has at least $2k$ vertices. Equality
holds exactly for the affine images of $C^\Delta_k$, and these have property (F). Hence,
under each reading above, the minimal number of vertices of a $k$-dimensional
cs-neighborly polytope is $2k$, attained exactly by the affine cross-polytopes.
\end{proposition}
\begin{proof}
Translate the center to $0$, so $V=-V$. Since $0=\frac12(v+(-v))$ lies in the
affine hull, the affine hull equals the linear span of $V$, which is therefore $\mathbb R^k$.
Choose a basis $v_1,\dots,v_k\in V$. The $2k$ points $\pm v_i$ lie in $V$ and are distinct:
$v_i\ne0$, and $v_i=\pm v_j$ for $i\ne j$ contradicts independence. If $|V|=2k$, then
$V=\{\pm v_i\}=A V_k$ with $A$ the invertible linear map $e_i\mapsto v_i$. An invertible
affine map preserves faces, central symmetry and antipodality, so $A C^\Delta_k$ has
property (F). Conversely, every affine image of $C^\Delta_k$ has $2k$ vertices.
\end{proof}

\begin{remark}[A ``maximal'' reading]
If ``minimal'' were a slip for ``maximal'', the claim would still fail for the $k$-dimensional
reading. Under the literal definition there is no maximum (centrally symmetric $2m$-gons
exist for every $m$). Under Gr\"unbaum's cs-$2$-neighborliness, Linial and Novik~\cite{LN}
proved that a cs-$2$-neighborly $d$-polytope has at most $2^d$ vertices, so $2^{k+1}$ is not
attained. This remark is not formalized.
\end{remark}

\paragraph{Hanner polytopes.}
Hanner polytopes are generated from the segment $[-1,1]$ by Cartesian products
$P\times Q$ and free sums $P\oplus Q=\operatorname{conv}(P\times\{0\}\cup\{0\}\times Q)$.
For centrally symmetric $P,Q$ with $0$ in the interior,
$f_0(P\times Q)=f_0(P)f_0(Q)$ and $f_0(P\oplus Q)=f_0(P)+f_0(Q)$, and dimensions add.
By induction, a $d$-dimensional Hanner polytope satisfies $2d\le f_0\le 2^d$. The base
case is the segment, $f_0=2$. For the product step, $d_1+d_2\le 2d_1d_2$ and $2^{d_1}2^{d_2}=2^{d_1+d_2}$.
For the free-sum step, $x+y\le xy$ for $x,y\ge2$ gives $2^{d_1}+2^{d_2}\le2^{d_1+d_2}$.
So no Hanner polytope has $2^{d+1}$ vertices, and the conjecture's two clauses cannot both
hold if ``extremal bodies'' means the bodies attaining the minimum. The cube ($2^d$
vertices) and the cross-polytope ($2d$ vertices) are Hanner. By
Proposition~\ref{prop:true}, the true minimizers are affine cross-polytopes, which are
Hanner up to linear equivalence. So the corrected statement ``minimum $2k$, attained
by (affine) Hanner polytopes'' is true. (The conjecture may be a garbled form
of Kalai's $3^d$ conjecture on the total number of faces of cs polytopes, where Hanner
polytopes are the conjectured extremal bodies. That is a different statement.)

\section{Lean formalization}
The project \texttt{lean4/} (Lean 4.19.0, core library only, no Mathlib) defines all
objects from scratch.
\begin{itemize}
\item Points of $\mathbb Z^k$ are lists of $k$ integers, with the scalar product \texttt{dot}.
A polytope is given by its vertex list \texttt{V}.
  \begin{itemize}
  \item \texttt{IsVertex V v}: $v$ is the unique maximizer over $V$ of some functional.
  \item \texttt{ConvexPosition V}: no repetitions, and every point is a vertex, so the
  polytope has \texttt{V.length} vertices.
  \item \texttt{IsFace V S}: $S$ is exactly the set of maximizers of some functional.
  \item \texttt{CentSymm V}: $V=-V$.
  \end{itemize}
\item \texttt{FullDim k V} is a dimension certificate: points $a,b_1,\dots,b_k\in V$ and
functionals $f_i$ with $f_i(b_j-a)\ne0\iff i=j$. \texttt{fullDim\_linIndep} and
\texttt{fullDim\_sound} prove that the $b_j-a$ are then linearly independent (the same
argument works over $\mathbb R$). Together with all points lying in $\mathbb Z^k$, this
gives \texttt{IsPolytopeOfDim k V}.
\item The three readings are \texttt{StatementCsNeighborly} (the literal definition),
\texttt{CsNeighborlyUpTo j} (Gr\"unbaum's cs-$j$-neighborliness) and \texttt{CsNeighborlyFull}
(property (F)). \texttt{full\_upTo} shows that (F) implies the second for every $j$, and
\texttt{statement\_def\_of\_cs} shows that every nonempty cs polytope satisfies the first.
\item \texttt{MinVerticesIs N k n} says that the minimum is attained and is a lower bound;
\texttt{MinClause N} is the clause $\forall k\ge1$, \texttt{MinVerticesIs N k (2\^{}(k+1))}.
\item \texttt{crossFacts\_1} to \texttt{crossFacts\_4} and \texttt{crossFacts\_6} check, by \texttt{decide} for the
generated list \texttt{cross k}:
  \begin{itemize}
  \item length $2k$, no repetitions, central symmetry;
  \item the vertex property (witness $c=v$) and the dimension certificate;
  \item for every subset $S$ (every sublist), that if $S$ is nonempty and antipodal-free,
  then the maximizers of $\sum S$ are exactly $S$.
  \end{itemize}
  \texttt{full\_of\_sublists} transfers the last check from sublists to arbitrary
  subsets.
\item Main theorems:
  \begin{itemize}
  \item \texttt{conjecture\_00000001564\_false : \textlnot\ MinClause StatementCsNeighborly};
  \item \texttt{conjecture\_00000001564\_false\_grunbaum (j)}, for every $j$;
  \item \texttt{conjecture\_00000001564\_false\_full};
  \item \texttt{conjecture\_00000001564\_false\_any}, for any notion satisfied by the octahedron;
  \item \texttt{conjecture\_00000001564\_false\_param}, the reading in Remark~\ref{rem:readings}(2);
  \item \texttt{conjecture\_00000001564\_false\_2k}: the Chinese $2k$-dimensional reading
  (Section~\ref{sec:cn}), via \texttt{crossFacts\_6}, with four instances (suffixes
  \texttt{statement}, \texttt{full}, \texttt{grunbaum}, \texttt{cs3});
  \item \texttt{fails\_each\_k}: failure at each $k=1,2,3,4$;
  \item \texttt{conjecture\_00000001564\_conjunction\_false}: the conjunction with any Hanner clause is false.
  \end{itemize}
\item Non-vacuity:
  \begin{itemize}
  \item \texttt{true\_min\_k1} proves a true instance, \texttt{MinVerticesIs \_ 1 2}, including its lower bound.
  \item A centrally symmetric hexagon is a 2-polytope satisfying the literal definition
  but not cs-$2$-neighborliness (\texttt{hexagon\_not\_csNeighborly}).
  \item A set with an interior point is not in convex position.
  \end{itemize}
\item \texttt{HannerCount} records the (dimension, vertex number) pairs generated by
segments, products and free sums. \texttt{hanner\_bounds} proves $2d\le v\le2^d$, and
\texttt{no\_hanner\_with\_2pow\_succ} proves $v\ne2^{d+1}$.
\end{itemize}
There is no \texttt{sorry}, \texttt{native\_decide} or added axiom. \texttt{\#print axioms}
reports only \texttt{propext}, \texttt{Quot.sound} and (for the generic linear-independence
lemma) \texttt{Classical.choice}.

\paragraph{Scope.}
\begin{itemize}
\item Lean checks the cross-polytope for $k=1,2,3,4,6$; one $k$ suffices to refute the clause.
The case of general $k$ is Theorems~\ref{thm:cross} and~\ref{thm:main} above, and
\texttt{verify.py} checks $k\le6$.
\item The general lower bound $2k$ of Proposition~\ref{prop:true} is proved here. It is
formalized only for $k=1$, and \texttt{verify.py} checks it for lattice sets in $\{-1,0,1\}^k$, $k=2,3$.
\item For Hanner polytopes, only the vertex-count recursion is formalized. The facts
$f_0(P\times Q)=f_0(P)f_0(Q)$ and $f_0(P\oplus Q)=f_0(P)+f_0(Q)$ are standard, and
\texttt{verify.py} confirms them on explicit point sets.
\end{itemize}
The integer functionals used as witnesses are real functionals, so every Lean predicate
implies its real counterpart. In particular, $\operatorname{conv}(\texttt{cross 3})$ is a
genuine $3$-dimensional cs polytope in $\mathbb R^3$ with $6$ vertices and property (F).

\section{Independent check (\texttt{verify.py})}
The Python script uses only the standard library and a method different from the Lean proof.
\begin{itemize}
\item For $k\le6$, it enumerates all functionals $c\in\{-1,0,1\}^k\setminus\{0\}$, and for
$k\le3$ also all $c\in\{-2,\dots,2\}^k\setminus\{0\}$. It collects every maximizer set and
confirms that this family is exactly the family of antipodal-free vertex subsets ($2,8,26,80,242,728$
for $k=1,\dots,6$). It also computes the dimension as an exact rank.
\item It checks every centrally symmetric point set in $\{-1,0,1\}^k$ ($k=2,3$) that
spans $\mathbb R^k$: each gives a polytope with at least $2k$ vertices.
\item It builds all $1,2,6,22,90,394$ explicit Hanner polytopes of dimension $1,\dots,6$
and confirms that their vertex numbers lie in $[2d,2^d]$ and never equal $2^{d+1}$.
\end{itemize}

\section{Reproduction}
\begin{verbatim}
cd lean4 && lake build        # Lean 4.19.0, core only
cd .. && python3 verify.py    # Python 3 standard library
pdflatex report.tex && pdflatex report.tex
\end{verbatim}

\begin{thebibliography}{9}
\bibitem{Gr} B.~Gr\"unbaum, \emph{Convex Polytopes}, 2nd ed., Springer GTM 221, 2003
(chapter on centrally symmetric polytopes).
\bibitem{Ha} O.~Hanner, Intersections of translates of convex bodies,
\emph{Math.\ Scand.} 4 (1956), 65--87.
\bibitem{LN} N.~Linial and I.~Novik, How neighborly can a centrally symmetric polytope be?,
\emph{Discrete Comput.\ Geom.} 36 (2006), 273--281.
\bibitem{Ka} G.~Kalai, The number of faces of centrally-symmetric polytopes,
\emph{Graphs Combin.} 5 (1989), 389--391.
\end{thebibliography}
\end{document}
